\documentclass[10pt,a4paper]{amsart}
\usepackage[margin=19mm]{geometry}
\usepackage{amsmath,amssymb,mathtools}
\usepackage[scr=rsfs,cal=euler]{mathalfa}
\usepackage{xfrac}
\usepackage[unicode,hidelinks,bookmarks=false]{hyperref}
\newcommand{\ie}{i.e.\ }
\newcommand{\cf}{cf.\ }
\newcommand{\resp}{respectively\ }
\newcommand{\Subsection}{\subsection}
\providecommand{\nobreakdash}{\nobreak\hbox{-}}
\setlength{\emergencystretch}{2em}
\multlinegap0pt
\usepackage{bm}
\usepackage{bbm}
\usepackage{dsfont}
\usepackage{xspace}
\usepackage{cancel}
\usepackage{mathtools}
\usepackage{amsthm}
\makeatletter
\@ifundefined{theorem}{\newtheorem{theorem}{Theorem}}{}
\@ifundefined{lemma}{\newtheorem{lemma}[theorem]{Lemma}}{}
\makeatother
\title[Local square mean in the hyperbolic circle problem and sums ]{Local square mean in the hyperbolic circle problem and sums of Sali\'e sums}
\author{Andr\'as Bir\'o}
\date{Source: arXiv:2604.11205v1 (CC BY 4.0)}
\begin{document}
\maketitle

\noindent\emph{Source and licence.} Local square mean in the hyperbolic circle problem and sums of Sali\'e sums. Version arXiv:2604.11205v1 (CC BY 4.0), distributed under the Creative Commons Attribution 4.0 International licence, as recorded in the arXiv licence metadata for this exact version and shown on the abstract page https://arxiv.org/abs/2604.11205v1. Full rights notice: licenses/arxiv-2604.11205-CC-BY-4.0.txt.\par\smallskip

\noindent\textit{Editorial scope.} Counted unit: the Lemma whose source label is \texttt{lem:12,9} in the article (source0.tex lines 1509--1527, source label \texttt{lem:12,9}), delivered together with the article's own complete original proof of it (source0.tex lines 1528--1594, one \texttt{proof} environment of 67 lines). No mathematics is written, repaired or re-derived: statement and proof are the article's own text line for line. Required context retained uncounted: 679 (lines 1284--1284); 680 (lines 1290--1292); 95 (lines 1301--1304); 96 (lines 1305--1305). Every definition, display and label of those lines is retained unchanged. Omitted from this package and not redistributed: 1--1283, 1285--1289, 1293--1300, 1306--1508, 1595--2647. Nothing inside a retained excerpt is omitted or altered: every retained body is delivered exactly as the article's lines are written, so no omission receipt of any kind accompanies this package. Citations: the retained excerpts carry no citation command at all, so no bibliography entry of the article is transcribed and no reference is redistributed. Reference adaptation: none was needed and none was made; every reference inside the delivered unit resolves inside it, so no unresolved reference or citation is produced. Printed slips kept literal: no printed word, symbol, hypothesis or number of the counted statement or of its proof was altered. Layout and class: the article's own class is \texttt{amsart} with packages including amssymb; the wrapper is the lane's pinned \texttt{amsart} editorial wrapper at 10pt on A4, which restates the theorem-like environments the retained text needs and adds the article's own macro definitions that the retained text needs (none), with extra wrapper packages (bm, bbm, dsfont, xspace, cancel, mathtools). The delivered numbering is the unit's own. Assets: no retained span contains a figure environment, a graphics-inclusion command, a PDF-page inclusion, a table or a TikZ picture; no image file is copied into this package, the package declares no asset dependency and no image file is redistributed with it. Licence: version arXiv:2604.11205v1 of this article is distributed under the Creative Commons Attribution 4.0 International licence, as recorded in the arXiv licence metadata for this exact version and shown on the abstract page https://arxiv.org/abs/2604.11205v1. A full rights notice is delivered at licenses/arxiv-2604.11205-CC-BY-4.0.txt. Characters: every retained excerpt is pure ASCII, so the delivered engine, XeLaTeX with its default Latin Modern text font, reports no missing character.
\begin{equation}\gamma\left(G,R\right):=16RG\prod_{p|G}p,\label{679}\end{equation}

\begin{equation}\prod_{p|G}p^{\nu_p\left(C^2-\left(A^2-4\right)\left
(B^2-4\right)\right)}=G,\;\prod_{p|G}p^{\nu_p\left(A^2-4\right)}=
R,\label{680}\end{equation}

\begin{equation}t_1\equiv A\left({\rm m}{\rm o}{\rm d}\,\gamma\left
(G,R\right)\right),t_2\equiv B\left({\rm m}{\rm o}{\rm d}\,\gamma\left
(G,R\right)\right),f\equiv C\left({\rm m}{\rm o}{\rm d}\,\gamma\left
(G,R\right)\right),\label{95}\end{equation}

\begin{equation}p|2\gcd\left(t_1^2-4,f\right)\Leftrightarrow p|G,\label{96}\end{equation}

\begin{lemma}\label{lem:12,9} Let $G$ and $R$ be positive integers such that $
4|G$ and
$p|2R\Leftrightarrow p|G$, and let $\left(A,B,C\right)\in H_{G,R}$. Let $
t_1,t_2>2$ and $f$ be integers such that
$f^2>\left(t_1^2-4\right)\left(t_2^2-4\right)$, $f^2\le 4X^2$ and \eqref{95} and \eqref{96}
are true. Let us write $N:=\frac {f^2-\left(t_1^2-4\right)\left(t_
2^2-4\right)}G$. Then $N$ is
an integer, $\gcd\left(N,G\right)=1$ and we have that $\alpha_G\left
(t_1,t_2,f\right)$ equals the sum of
\[\sum_{1\le e^m\le 2eX}\sum_{D|N}\left(\frac {t_1^2-4}D\right)\phi_
0\left(\frac D{e^m}\right)\]
and
\[\kappa\sum_{2eX<e^m\le 4eX^2}\sum_{1\le e^r\le 2eX}\sum_{D|N}\left
(\frac {t_1^2-4}D\right)\phi_0\left(\frac D{e^r}\right)\frac 1{2\pi
i}\int_{\left(0\right)}\left(\frac {De^m}N\right)^{-s}\Phi_0\left
(s\right)ds,\]
where $m$ and $r$ run over integers, $\kappa :=\kappa_{G,R,A,B,C}$ depends only on $
G,R,A,B$ and $C$, and we have $\kappa_{G,R,A,B,C}\in \{-1,1\}$.
\end{lemma}

\begin{proof}
By \eqref{96}, \eqref{680}, \eqref{95} and \eqref{679} we have that $
p|2\gcd\left(t_1^2-4,f\right)\Leftrightarrow p|G$ and
\[\prod_{p|G}p^{\nu_p\left(f^2-\left(t_1^2-4\right)\left(t_2^2-4\right
)\right)}=G,\;\prod_{p|G}p^{\nu_p\left(t_1^2-4\right)}=R.\]
In particular, $N$ is an integer and $\gcd\left(N,G\right)=1$.

Then we have that if $D|f^2-\left(t_1^2-4\right)\left(t_2^2-4\right
)$ and $\gcd\left(D,G\right)=1$, then
\begin{equation}\left(\frac {t_1^2-4}D\right)=\left(\frac RD\right
)\left(\frac {\left(t_1^2-4\right)/R}D\right).\label{501}\end{equation}
Since $2|G$, we see that $\left(t_1^2-4\right)/R$ and $D$ are odd. By
quadratic reciprocity we have
\begin{equation}\left(\frac {\left(t_1^2-4\right)/R}D\right)=\left
(-1\right)^{\frac {D-1}2\frac {\frac {t_1^2-4}R-1}2}\left(\frac D{\left
(t_1^2-4\right)/R}\right),\label{502}\end{equation}
and
\begin{equation}\left(\frac D{\left(t_1^2-4\right)/R}\right)=\left
(\frac {\left(f^2-\left(t_1^2-4\right)\left(t_2^2-4\right)\right)
/D}{\left(t_1^2-4\right)/R}\right).\label{503}\end{equation}
To see this last statement we have to check that
$\gcd\left(f,\frac {t_1^2-4}R\right)=1$. But this is true, since a common prime
divisor would divide $G$, but $\gcd\left(G,\frac {t_1^2-4}R\right
)=1$.
We see that
\begin{equation}\left(\frac {\left(f^2-\left(t_1^2-4\right)\left(
t_2^2-4\right)\right)/D}{\left(t_1^2-4\right)/R}\right)=\left(\frac {
D^{\ast}}{\left(t_1^2-4\right)/R}\right)\left(\frac G{\left(t_1^2
-4\right)/R}\right),\label{504}\end{equation}
where
\begin{equation}D^{\ast}:=\frac {f^2-\left(t_1^2-4\right)\left(t_
2^2-4\right)}{DG}.\label{505}\end{equation}
Again by qauadratic reciprocity we get
\begin{equation}\left(\frac {D^{\ast}}{\left(t_1^2-4\right)/R}\right
)=\left(-1\right)^{\frac {D^{\ast}-1}2\frac {\frac {t_1^2-4}R-1}2}\left
(\frac {t_1^2-4}{D^{\ast}}\right)\left(\frac R{D^{\ast}}\right).\label{506}\end{equation}
By \eqref{501}, \eqref{502}, \eqref{503}, \eqref{504},
\eqref{506} and \eqref{505} we obtain
\begin{equation}\left(\frac {t_1^2-4}D\right)=\left(\frac {t_1^2-
4}{D^{\ast}}\right)\kappa\label{507}\end{equation}
with
\[\kappa :=\left(\frac R{\frac {f^2-\left(t_1^2-4\right)\left(t_2^
2-4\right)}G}\right)\left(-1\right)^{\frac {D+D^{\ast}-2}2\frac {\frac {
t_1^2-4}R-1}2}\left(\frac G{\left(t_1^2-4\right)/R}\right).\]
We have that $4|$$\left(D-1\right)\left(D^{\ast}-1\right)$, since $
D$ and $D^{\ast}$ are odd.
Therefore
\[\frac {D+D^{\ast}-2}2-\frac {DD^{\ast}-1}2\]
is even. Consequently, we get by \eqref{505} that $\kappa$ equals
\[\left(\frac R{\frac {f^2-\left(t_1^2-4\right)\left(t_2^2-4\right
)}G}\right)\left(-1\right)^{\frac {\frac {f^2-\left(t_1^2-4\right
)\left(t_2^2-4\right)}G-1}2\frac {\frac {t_1^2-4}R-1}2}\left(\frac
G{\left(t_1^2-4\right)/R}\right).\]
We see by \eqref{679} and \eqref{95} that $t_1$, $t_2$ and $f$
are fixed modulo $4GR$, hence $\kappa$ is determined by $G,R,A,B$
and $C$. Then by \eqref{505} and \eqref{507} we obtain
that $\alpha_G\left(t_1,t_2,f\right)$ equals the sum of
\[\sum_{1\le e^m\le 2eX}\sum_{D|N}\left(\frac {t_1^2-4}D\right)\phi_
0\left(\frac D{e^m}\right)\]
and
\[\kappa\sum_{2eX<e^m\le 4eX^2}\sum_{D|N}\left(\frac {t_1^2-4}D\right
)\phi_0\left(\frac {De^m}N\right).\]
Here we must have $\frac {De^m}N\le e$, but $N\le 4X^2,$ $e^m>2eX$,
hence $1\le D\le 2X$. So we may insert the factor
$\sum_{1\le e^r\le 2eX}\phi_0\left(\frac D{e^r}\right)$. The lemma follows by Mellin
inversion.
\end{proof}

\end{document}
